% !TEX root = ./main.tex

\begin{center}
\textbf{Instruções para elaboração do Trabalho Final (artigo/ensaio)}
\end{center}
\bigskip

\noindent\textbf{O Relatório Final deve ser produzido observando-se os seguintes elementos:}
\begin{taglist}
  \item Introdução (contemplando objetivos, justificativa, problemática abordada);
  \item Materiais e método;
  \item Resultados e Discussão;
  \item Conclusões/Considerações finais;
  \item Referências.
\end{taglist}

\bigskip
\noindent\textbf{O relatório deve ser produzido conforme normas da ABNT NBR 14724/2024 para normas de trabalhos acadêmicos, ABNT NBR 10520:2023 para citações e ABNT NBR 6023:2018 para referências, observando:}
\begin{taglist}
  \item Quantidade de palavras: entre 4 e 6 mil palavras ou de 10 a 15 páginas (sem contar as referências bibliográficas)
  \item Margens: Esquerda e Superior: 3cm, Direita e Inferior: 2cm;
  \item Fonte: Times New Roman 12;
  \item Espaçamento entre linhas 1,5
  \item Citações longas (acima de 3 linhas): Tamanho da Fonte 11; Espaçamento simples; Recuo na margem esquerda: 4cm.
\end{taglist}

\bigskip
\instrucao{Redigir este Trabalho Final (artigo/ensaio) em arquivo separado do Relatório Técnico Final;}
\instrucao{Pesquisas financiadas pelo CNPq, sempre que possível, apresentar relação com alguma das áreas prioritárias do CNPQ conforme estabelecido na portaria MCTI nº 1.122/2020, com texto alterado pela portaria MCTI nº 1.329/2020;}
\instrucao{Salvar em PDF para enviar como documento suplementar no sistema do IC\&ITI;}
\instrucao{Caso o discente tenha submetido os resultados da pesquisa de IC ou ITI em periódicos ou eventos científicos (como trabalho completo), após o mês de maio de 2026, poderá substituir o envio deste artigo/ensaio (modelo PRPPG) pelo submetido à revista científica ou ao evento científico.}
\instrucao{Este mesmo artigo/ensaio poderá ser enviado como trabalho completo do SIPEC, caso seja de interesse do autor ou autora;}

\begin{center}
\colorbox{docred}{\textbf{\textcolor{black}{APAGAR ESTAS INSTRUÇÕES AO FINAL}}}
\end{center}

\bigskip

% =====================================================================
% INTRODUÇÃO
% =====================================================================
\seguetitulo{INTRODUÇÃO}{(em caixa alta, negrito, sem número indicativo de seção, sem ponto, fonte Times New Roman, tamanho da fonte 12, alinhado à esquerda, espaçamento entre linhas 1,5. Para iniciar o texto, deixa-se uma linha em branco)}

\medskip
A introdução deve apresentar os elementos constitutivos do problema e da problemática de pesquisa, além dos objetivos a serem trabalhados no decorrer do texto. Pode ainda conter o recorte teórico e temporal do objeto de investigação além da fundamentação teórica que embasou a pesquisa de Iniciação Científica.

Os parágrafos serão redigidos em Times New Roman, tamanho 12, alinhamento justificado com recuo de um TAB (1,25) na primeira linha e as margens devem ter a seguinte configuração: Esquerda e Superior: 3cm, Direita e Inferior: 2cm.

\bigskip
% =====================================================================
% MATERIAIS E MÉTODOS
% =====================================================================
\seguetitulo{MATERIAIS E MÉTODOS}{(em caixa alta, negrito, sem número indicativo de seção, sem ponto, fonte Times New Roman, tamanho da fonte 12, alinhado à esquerda, espaçamento entre linhas 1,5. Para iniciar o texto, deixa-se uma linha em branco)}

\medskip
Neste momento, apresentar os materiais utilizados como fontes e embasamento teórico, além do método e metodologias utilizados durante o período de vigência da pesquisa científica.

\bigskip
\noindent\hspace*{1.25cm}\textbf{Observações:}

\medskip
\begin{adjustwidth}{4cm}{0cm}
{\singlespacing\fontsize{10}{12}\selectfont
As citações diretas que contenham mais do que três linhas, deverão ser apresentadas em conformidade com a ABNT (NBR 10520/2023) em parágrafo especial, separado do texto por uma linha em branco antes e depois da citação, com recuo de 4 cm e tamanho da fonte 10, espaçamento entre linhas simples. (Sobrenome, ano, p.).
\par}
\end{adjustwidth}

\bigskip
As citações diretas, de até três linhas, serão integradas no texto e colocadas entre aspas duplas e toda vez que houver a ocorrência de supressão de texto antes, durante ou após a citação, esta deve ser indicada da seguinte forma: ``[...] nonono nonono nonono nonono [...] nonono nonono nonono nonono nonono nonono nonono [...]'' (Sobrenome, ano, p.). \textbf{Atenção}, nunca termine uma seção, ou subseção com citação.

\bigskip
% =====================================================================
% RESULTADOS E DISCUSSÕES
% =====================================================================
\seguetitulo{RESULTADOS E DISCUSSÕES}{(em caixa alta, negrito, sem número indicativo de seção, sem ponto, fonte Times New Roman, tamanho da fonte 12, alinhado à esquerda, espaçamento entre linhas 1,5)}

\medskip
Caso necessário, neste item, poderá haver seção e subseções. A quantidade de seções e subseções do trabalho completo pode variar de acordo com a necessidade, ficando a critério do autor ou autora. Além disso, todas as seções devem conter um texto relacionado a elas.

\bigskip
\seguetitulo{TÍTULO DA SEÇÃO}{(O título das seções será digitado em negrito, com letras maiúsculas, sem recuo, sem ponto e sem numeração. As subseções serão digitadas em negrito. Entre o final de uma seção ou subseção e a parte subsequente, deixa-se uma linha em branco)}

\bigskip
\noindent\textbf{Título da Subseção} {\color{docred}(primeira letra em maiúscula, negrito, sem número indicativo de seção, sem ponto, fonte Times New Roman, tamanho da fonte 12, alinhado à esquerda, espaçamento entre linhas 1,5. Entre o final de uma seção ou subseção e a parte subsequente, deixa-se uma linha em branco)}

\bigskip
As ilustrações (quadros, figuras, fotos, gráficos, etc.) devem localizar-se o mais próximo possível do texto a que se referem e apresentar uma numeração sequencial em algarismos arábicos de acordo com a ordem de ocorrência no texto.

\bigskip
\noindent\hspace*{1.25cm}\textbf{Imagem 1 -- Título da Imagem}

\medskip
\begin{center}
\fbox{\parbox[c][2.8cm]{9cm}{\centering IMAGEM}}
\end{center}

\medskip
\noindent Fonte: Fonte da Imagem (Autor, ano, p).

\bigskip
As tabelas ou quadros também devem ter um número sequencial em algarismos arábicos, inscrito na parte superior, a esquerda da página, precedida da palavra ``Tabela'' ou ``Quadro'' quando for o caso (sempre que o artigo apresentar duas ou mais tabelas ou quadros). Além disso, cada uma deve possuir um título por extenso, inscrito no topo. A referência, caso seja uma citação de outra obra, deve ser colocada imediatamente abaixo da tabela ou quadro, precedida da palavra Fonte. Veja o exemplo abaixo:

\bigskip
\noindent\hspace*{1.25cm}\textbf{Quadro 1 -- Título do Quadro}

\medskip
\begin{center}
\begin{tabular}{|p{3cm}|p{3cm}|p{3cm}|p{3cm}|}
\hline
\multicolumn{4}{|c|}{Nonono} \\
\hline
nonono & Nonono & nonono & Nonono \\
\hline
 & & & \\
\hline
 & & & \\
\hline
\end{tabular}
\end{center}

\medskip
\noindent Fonte: Fonte do Quadro (Autor, ano, p).

\bigskip
\noindent\hspace*{1.25cm}\textbf{Tabela 1 -- Título da Tabela}

\medskip
\begin{center}
\begin{tabular}{cccc}
\toprule
nonono & Nonono & nonono & Nonono \\
\midrule
 & & & \\
\midrule
 & & & \\
\bottomrule
\end{tabular}
\end{center}

\medskip
\noindent Fonte: Fonte da Tabela (Autor, ano, p).

\bigskip
% =====================================================================
% CONSIDERAÇÕES FINAIS
% =====================================================================
\seguetitulo{CONSIDERAÇÕES FINAIS}{(em caixa alta, negrito, sem número indicativo de seção, sem ponto, fonte Times New Roman, tamanho da fonte 12, alinhado à esquerda, espaçamento entre linhas 1,5. Para iniciar o texto, deixa-se uma linha em branco)}

\medskip
Não é indicado utilizar citação nas considerações finais, assim como assuntos que não foram tratados no corpo do texto.

\bigskip
% =====================================================================
% REFERÊNCIAS BIBLIOGRÁFICAS
% =====================================================================
\seguetitulo{REFERÊNCIAS BIBLIOGRÁFICAS}{(em caixa alta, negrito, sem número indicativo de seção sem ponto, fonte Times New Roman, tamanho da fonte 12, alinhado à esquerda, espaçamento entre linhas 1,5. Para iniciar o texto, deixa-se uma linha em branco)}

\medskip
\begingroup\singlespacing
\noindent\textcolor{docblue}{\textbf{ATENÇÃO}}: as referências devem ser apresentadas \textcolor{docblue}{\textbf{em ordem alfabética}} de sobrenome do autor-data e alinhadas à margem esquerda, espaçamento entre linha de 1,0 em conformidade com a ABNT (NBR 6023/2018), seguindo, em linhas gerais, as orientações abaixo:
\par

\bigskip
\noindent\textcolor{docblue}{\textbf{[Obras no todo]}}\\
SOBRENOME, Nome do autor. \textbf{Título da obra.} Cidade: Editora, Ano.

\bigskip
\noindent\textcolor{docblue}{\textbf{[Parte de obra com autor específico]}}\\
SOBRENOME, Nome do autor do capítulo, da seção ou do artigo. Título do capítulo, seção ou artigo. \textit{In}: SOBRENOME, Nome do autor/organizador (org.). \textbf{Título da obra.} Cidade: Editora, p. x-y (número das páginas do capítulo ou artigo), Ano.

\bigskip
\noindent\textcolor{docblue}{\textbf{[Monografia no todo]}}\\
SOBRENOME, Nome do autor. \textbf{Título da monografia.}  Monografia (Departamento no qual a monografia foi defendida) -- Nome da Universidade, Cidade, Ano.

\bigskip
\noindent\textcolor{docblue}{\textbf{[Artigo de uma revista com autor definido]}}\\
SOBRENOME, Nome do autor. Título do artigo. \textbf{Nome da revista}, Cidade, volume, número da revista, p. x-y (número das páginas de início e fim do artigo), ano.

\bigskip
\noindent\textcolor{docblue}{\textbf{[Artigo de um jornal com autor definido]}}\\
SOBRENOME, Nome do autor. Título do artigo ou da matéria, subtítulo (se houver). \textbf{Título do Periódico}, Local de publicação, numeração do ano e/ou volume, número e/ou edição, tomo (se houver), páginas inicial e final, e data ou período de publicação.

\bigskip
\noindent\textcolor{docblue}{\textbf{[Dissertação/Tese no todo]}}\\
SOBRENOME, Nome do autor. \textbf{Título da tese.} Tese (Doutorado em -- área de estudo) -- Universidade em que a tese foi defendida, Local, ano.

\bigskip
\noindent\textcolor{docblue}{\textbf{[Texto em meio eletrônico]}}\\
SOBRENOME, Nome do Autor. \textbf{Título}. Disponível em: http://www.endereço completo do site. Acesso em: dia, mês (abreviado), ano.

\bigskip
\noindent\textcolor{docblue}{\textbf{[Trabalho apresentado em evento]}}\\
SOBRENOME, Nome do Autor. Título. \textit{In:} NOME DO EVENTO, número do evento, ano, cidade. \textbf{Anais} [...] (nome dos anais). Cidade: Instituição organizadora do evento, ano. p. x-y (páginas da publicação nos anais do evento). *

\medskip
\noindent * se os anais estiverem online, acrescentar os elementos de publicação em meio eletrônico: Disponível em: http://www.endereço completo do site. Acesso em: dia, mês (abreviado), ano.

\bigskip
\noindent\textcolor{docblue}{\textbf{[Legislação]}}\\
JURISDIÇÃO OU CABEÇALHO DA ENTIDADE. [Epígrafe (ano)]. \textbf{Ementa transcrita conforme publicada}. Local de Publicação: Editora, ano.

\medskip
\noindent Exemplo: RIO GRANDE DO SUL. [Constituição (1989)]. \textbf{Constituição do Estado do Rio Grande do Sul}. 4. ed. atual. Porto Alegre: Assembléia Legislativa do Estado do Rio Grande do Sul, 1995.

\bigskip
\noindent\textcolor{docblue}{\textbf{[Obra com vários autores]}}\\
SOBRENOME, Nome do Primeiro Autor \textit{et al}. Título da obra. Cidade: Editora, ano.

\medskip
\noindent Obs.: Em casos específicos (projetos de pesquisa científica, indicação de produção científica em relatórios para órgãos de financiamento etc.), nos quais a menção dos nomes for indispensável para certificar a autoria, é facultado indicar todos os nomes.

\bigskip
\noindent\textcolor{docblue}{\textbf{[Publicação Audiovisual -- filmes, vídeos, entre outros]}}\\
NOME do vídeo. Direção: Nome do Diretor. Produção: Nome do(s) Produtor(es). Cidade: Editora ou Produtora, ano. 1 DVD (indicar o suporte e em unidades físicas), (indicar dentro dos parênteses a duração em minutos).

\bigskip
\noindent\textcolor{docblue}{\textbf{[Documentos iconográficos -- pinturas, gravuras, ilustração, fotografia, desenho técnico, entre outros]}}\\
SOBRENOME, Nome do Autor. \textbf{Título da obra.} ano. 1 fotografia. (especificação do suporte de apresentação da obra, pintura, fotografia, gravura...)

\bigskip
\noindent\textcolor{docblue}{\textbf{[Documento sonoro no todo]}}\\
TÍTULO da obra. Compositor: Nome do Compositor. Intérprete: Nome do Intérprete. Cidade: Gravadora, ano. 1 CD (especificação do tipo de suporte e quantidade) (indicar dentro dos parênteses a duração em minutos.
\endgroup